\clearpage

% --------------------------------------------------------------------------------------------

\section {Summary}
\label{sec:summary}

In this technical note,
we have presented the results of a 3-flavour $\nu_{\mu}$-disappearance analysis,
performed by the Valencia-Lancaster-Oxford-Rutherford (VaLOR) group,
on the combined Run1+2+3+4a (6.3933$\times$10$^{20}$ POT) dataset.
This analysis is an update of the work presented in
Refs. \cite{OurRun1DisappAnalysis}, \cite{OurRun12DisappAnalysis}, 
\cite{OurRun12DisappAnalysis3Flavour} and
\cite{OurRun123DisappAnalysis3Flavour}
with improved methods and an enlarged dataset.

This analysis predicts 421.76 $\pm$ 23.33 (syst)
1 $\mu$-like ring events in SuperK in the absence of any oscillation,
but only 115 were observed. The observed deficit has a strong energy dependence;
the ratio of observed to expected, under the no oscillation hypothesis, is
$\sim$27\% $<$ 0.5 GeV, 
$\sim$10\% between 0.5 and 1 GeV and 
$\sim$73\% $>$ 1 GeV.

A $\nu_{\mu}$-disappearance analysis was performed 
in a framework of 3-flavour oscillations including matter effects in constant-density matter.
A total of 44 energy scale, neutrino flux, neutrino cross section,
efficiency, initial-state nuclear environment and final-state rescattering,
and oscillation parameter systematics were taken into account. 
The effect of these systematic uncertainties on the expected number of 1 $\mu$-like ring events,
on their reconstructed energy spectrum and on our oscillation measurement were presented in detail.

A 3-flavour $\nu_{\mu}$-disappearance fit to
the observed reconstructed energy spectrum of 1 $\mu$-like ring events of the combined Run 1+2+3+4 dataset,
allowing all nuisance parameters to float
% and with $\sin^{2}\theta_{23}$ constrained to the physical region,
yields
$|\Delta m^{2}_{32}|$ = 2.463$\times$10$^{-3}$ eV$^{2}$/c$^{4}$  and
$\sin^{2}\theta_{23}$ = 0.514
as the best-fit values (p-value = 0.88).
The fit assumed $\delta_{CP}$=0 and the normal mass hierarchy.
A fit assuming the inverted mass hierarchy yields
$|\Delta m^{2}_{32}|$ = 2.530$\times$10$^{-3}$ eV$^{2}$/c$^{4}$  and
$\sin^{2}\theta_{23}$ = 0.511
as the best-fit values (p-value = 0.89).

Confidence regions obtained using both the constant-$\Delta\chi^2$ method and the 
method of Feldman and Cousins were presented.
Using the profile $\Delta\chi^2$ distributions obtained from the fits,
the 90\%(68\%) CL allowed region from this analysis can be summarized as
2.30$\times$10$^{-3}$ (2.35$\times$10$^{-3}$) eV$^{2}$/c$^{4}$ 
$< |\Delta m^{2}_{32}| < 
$ 2.63$\times$10$^{-3}$ (2.57$\times$10$^{-3}$) eV$^{2}$/c$^{4}$ 
and 
0.433 (0.456) $< \sin^{2}\theta_{23} <$ 0.594 (0.572)
for the normal mass hierarchy,
and
2.37$\times$10$^{-3}$ (2.42$\times$10$^{-3}$) eV$^{2}$/c$^{4}$ 
$< |\Delta m^{2}_{32}| < 
$ 2.70$\times$10$^{-3}$ (2.64$\times$10$^{-3}$) eV$^{2}$/c$^{4}$ 
and 
0.431 (0.454) $< \sin^{2}\theta_{23} <$ 0.591 (0.569)
for the inverted mass hierarchy.

Numerous comparisons were made of the best-fit values and spectra 
for different running periods.
Comparisons were presented of the confidence regions obtained
by using different construction methods and by fits to different datasets.
Finally, comparisons of our contours with
recent measurements from MINOS and SuperK
% and the expected T2K sensitivity 
% for 6.393$\times$10$^{20}$ POT
were also made. 

